\documentclass[12pt,a4paper]{article}
\usepackage[margin=2cm]{geometry}
\usepackage{amsmath,amssymb,amsthm,mathtools}
\usepackage{enumitem}
\usepackage{algorithm}
\usepackage{algpseudocode}
\usepackage{microtype}
\usepackage[hidelinks]{hyperref}

\setlist{nosep}
\setlength{\parindent}{0pt}
\setlength{\parskip}{0.4em}

\newtheorem{theorem}{Theorem}[section]
\newtheorem{lemma}[theorem]{Lemma}
\newtheorem{proposition}[theorem]{Proposition}
\newtheorem{corollary}[theorem]{Corollary}

\theoremstyle{definition}
\newtheorem{definition}[theorem]{Definition}
\newtheorem{example}[theorem]{Example}
\newtheorem{problem}{Problem}[section]

\newtheorem*{remark}{Remark}

\newenvironment{solution}{\begin{proof}[Solution]}{\end{proof}}

\DeclareMathOperator{\OPT}{OPT}
\DeclareMathOperator{\poly}{poly}

\DeclareMathOperator{\ppi}{ppi}


\title{Placing fire trucks}
\author{Klevis Imeri}
\date{}

\begin{document}
\maketitle

First let us formalize what problem we are solving.

The input to the problem:
\[
\mathcal{I} := \left\{
\begin{aligned}
N &:= \{\text{villages}\}, && |N|=n,\\
P[i] &:= \text{population of village }i, && P[i]\geq0,\quad |P|=n,\\
TT &:= \{(t,r)\mid t\text{ is an allowed team for truck }r\}, && |TT|=m,\\
V[(t,r)] &:= \{v\in N\mid (t,r)\text{ covers }v\}, && (t,r)\in TT.
\end{aligned}
\right\}.
\]

The output is a set of pairs $\mathcal{S}=\{(t,r),\ldots\}$.
\[
\operatorname{feasible}(\mathcal{S}):=
\mathcal{S}\subseteq TT\ \land\
\forall (t,r),(t',r')\in\mathcal{S}:\;
\bigl((t,r)\neq(t',r')\implies t\neq t'\land r\neq r'\bigr)
.
\]
The cost
\[
f(\mathcal{S}) =
\sum_{v\in\bigcup_{(t,r)\in\mathcal{S}} V[(t,r)]} P[v].
\]

Our goal is to design an algorithm $\mathcal{A}$, with
$\mathcal{S}_{\mathcal{A}}=\mathcal{A}(\mathcal{I})$, such that:
\begin{enumerate}
\item $\mathcal{A}$ terminates.
\item $\operatorname{feasible}(\mathcal{S}_{\mathcal{A}})$.
\item $\operatorname{time}_{\mathcal{A}}(\mathcal{I})\leq\poly(|\mathcal{I}|)$.
\item $\displaystyle f(\mathcal{S}_{\mathcal{A}})
\geq \frac{1}{3}\max_{\substack{\mathcal{S}\subseteq TT\\\operatorname{feasible}(\mathcal{S})}} f(\mathcal{S})
=\frac{f(\mathrm{OPT})}{3}$ ($\mathcal{A}$ is a $3$-approximation).
\end{enumerate}

Let us design the algorithm $\mathcal{A}$. The first intuition is that this problem
is similar to set cover from the lecture, with at least these two differences:
\begin{enumerate}
\item In minimum set cover, we must cover every element. Selecting all sets would
achieve this, but our goal of minimizing cost works against selecting everything.
Here, we want to cover as much population as possible, but the feasibility
constraint prevents us from selecting all pairs, each team and each truck can
be used at most once. Thus, feasibility means covering every element in minimum
set cover, whereas here it means selecting pairs without team or truck conflicts.
\item In minimum set cover, we minimize the total cost of the selected sets.
Here, we maximize the total population covered.
So, instead of a minimization problem, we have a maximization problem.
\end{enumerate}

We hope that a similar greedy algorithm, with a choice similar to the one in
minimum set cover, will suffice for this problem. Therefore, in each phase,
we choose a pair covering the largest population among the villages not yet covered,
$N'$. We call the population newly covered by a pair its gain
(similar to the $ppi$ in the lecture):
\[
\operatorname{gain}((t,r)):=\sum_{v\in V[(t,r)]\cap N'}P[v].
\]

With this choice, we proceed to write the greedy algorithm:
\begin{algorithm}[H]
\caption{$\textit{placing\_fire\_trucks}(\mathcal{I})$}
\begin{algorithmic}
\If{$TT=\varnothing$ or $N=\varnothing$}
    \State \Return $\varnothing$
\EndIf
\State $(t,r)\gets\text{any element of }\arg\max_{(t',r')\in TT}\operatorname{gain}((t',r'))$
\State $TT'\gets\{(t',r')\in TT\mid t'\neq t\land r'\neq r\}$
\State $N'\gets N\setminus V[(t,r)]$
\State $\mathcal{I}'\gets\{N',P,TT',V\}$
\State \Return $\{(t,r)\}\cup\Call{placing\_fire\_trucks}{\mathcal{I}'}$
\end{algorithmic}
\end{algorithm}

Now we have to prove all four goals above for our algorithm.
\begin{enumerate}
\item Termination. If the algorithm does not return immediately, then
$TT\neq\varnothing$, so a maximizing pair $(t,r)$ exists. This pair is removed:
\[
|TT'|<|TT|\leq m,\qquad |N'|\leq|N|\leq n.
\]
By the well-foundedness of the nonnegative integers, $|TT|$ cannot decrease
strictly forever. Thus, after at most $m$ selections, the algorithm terminates.
\item Feasibility.
\begin{enumerate}
\item The first selected pair belongs to $TT$, and each subsequent selected pair
belongs to the remaining $TT'\subseteq TT$.
Therefore, $\mathcal{S}_{\mathcal{A}}\subseteq TT$.
\item Whenever we select $(t,r)$, we remove every pair $(t',r')$ with
$t'=t$ or $r'=r$. Thus, no pair selected in a subsequent iteration can share
a team or truck with $(t,r)$. Hence,
\[
\forall (t,r),(t',r')\in\mathcal{S}_{\mathcal{A}}:\quad
(t,r)\neq(t',r')\implies t\neq t'\land r\neq r'.
\]
\end{enumerate}
Both conditions hold throughout each iteration, so
$\operatorname{feasible}(\mathcal{S}_{\mathcal{A}})$.
\item Running time. Checking the stopping condition takes constant time. We compute the
gain of at most $m$ pairs, each covering at most $n$ villages, so this takes
$O(nm)$ time. Removing conflicting pairs takes $O(m)$ time, and removing covered villages takes $O(n)$ time.
Preparing the recursive call and combining its result take at most $O(n+m)$
additional time, using shared references for $P$ and $V$. Thus, with at most
$m$ iterations (each iteration removes at least one pair),
\[
T_i=O(nm),\qquad
T=\sum_{i=1}^{\text{number of iterations}}T_i
\leq m\cdot O(nm)
=O(nm^2)\subseteq O(n^2m^2)
\subseteq\poly(|\mathcal{I}|).
\]
\item $3$-approximation. Now we have to prove the important part. A usual trick
is to start with the following comparison in the subproblem of each iteration:

In the minimum case:
\[
\text{minimum over all available choices}
\leq
\text{minimum over the available choices in OPT}.
\]
In the maximum case:
\[
\text{maximum over all available choices}
\geq
\text{maximum over the available choices in OPT}.
\]

In minimum set cover, we used
$\min_{S_j\in\mathcal{S}'}ppi(S_j)\leq\min_{O_j\in\mathrm{OPT}'}ppi(O_j)$,
where $\mathcal{S}'$ and $\mathrm{OPT}'$ denote the choices remaining in the
current subproblem.
Here, we maximize gain, so whenever $\mathrm{OPT}'\neq\varnothing$, our comparison is
\begin{equation}
\label{eq:greedy-gain-comparison}
\max_{i\in TT'}\operatorname{gain}(i)\geq
\max_{o\in\mathrm{OPT}'}\operatorname{gain}(o),
\end{equation}
where $TT'$ is the set of currently available pairs and $\mathrm{OPT}'$
contains the pairs of a fixed optimal solution that are still available.

To simplify notation, for the pair $t$ selected by greedy, write:
\[
\operatorname{gain}_t(t):=\max_{i\in TT'_t}\operatorname{gain}_t(i)
\overset{\eqref{eq:greedy-gain-comparison}}{\geq}
\max_{o\in\mathrm{OPT}'_t}\operatorname{gain}_t(o)
\geq\operatorname{gain}_t(o),
\qquad\forall o\in\mathrm{OPT}'_t.
\]

Now we connect the comparison above to what happens in each iteration.
The chosen team and truck may belong to two different pairs in $\mathrm{OPT}'_t$.
Since each team and truck appears at most once in an optimal solution, selecting
$t$ removes at most two optimal pairs. Let $D_t$ be the set of these removed pairs.
If $\mathrm{OPT}'_t=\varnothing$, then $D_t=\varnothing$ and the per-iteration
bound below holds automatically.
By equation~\eqref{eq:greedy-gain-comparison}, each has gain at most that of $t$, so:
\[
|D_t|\leq2,\qquad
2\operatorname{gain}_t(t)\geq\sum_{o\in D_t}\operatorname{gain}_t(o).
\]
We use this form because summing the left-hand side over all iterations gives
exactly twice the total population (cost) covered by the greedy solution:
\[
2f(\mathcal{S}_{\mathcal{A}})
=2\sum_{t\in\mathcal{S}_{\mathcal{A}}}\operatorname{gain}_t(t)
\geq\sum_{o\in\mathrm{OPT}}\operatorname{gain}_{\tau(o)}(o).
\]
Here, $t\in\mathcal{S}_{\mathcal{A}}$ is a selected pair, and
$\operatorname{gain}_t$, $TT'_t$, and $\mathrm{OPT}'_t$ refer to the subproblem
in which $t$ is selected.

Now we have bounded $f(\mathcal{S}_{\mathcal{A}})$ by the expression on the right,
but how do we relate the right-hand side to $f(\mathrm{OPT})$?

Let us simplify the right-hand side. In each iteration, we remove the pairs
sharing the chosen team or truck. Each optimal pair is removed at most once.
If we finish with $TT=\varnothing$, every optimal pair has been removed, so:
\[
\bigcup_t D_t=\mathrm{OPT},\qquad D_t\cap D_s=\varnothing\quad(t\neq s).
\]
If we stop earlier because no uncovered villages remain, the remaining optimal
pairs have zero gain. Let $\tau(o)$ be the iteration in which $o$ is removed,
or the final state if $o$ remains. The gain $\operatorname{gain}_{\tau(o)}(o)$
is measured immediately before that iteration's pair is selected and its
covered villages are removed, or at termination if $o$ remains.
Thus, we can write the right-hand side as:
\[
\sum_{t\in\mathcal{S}_{\mathcal{A}}}\sum_{o\in D_t}\operatorname{gain}_t(o)
=\sum_{o\in\mathrm{OPT}}\operatorname{gain}_{\tau(o)}(o).
\]

To obtain $3f(\mathcal{S}_{\mathcal{A}})\geq f(\mathrm{OPT})$, it suffices to show:
\[
2f(\mathcal{S}_{\mathcal{A}})
\geq\sum_{o\in\mathrm{OPT}}\operatorname{gain}_{\tau(o)}(o)
\geq f(\mathrm{OPT})-f(\mathcal{S}_{\mathcal{A}}).
\]
The first inequality is already proved, so we only have to prove:
\[
\sum_{o\in\mathrm{OPT}}\operatorname{gain}_{\tau(o)}(o)
\geq f(\mathrm{OPT})-f(\mathcal{S}_{\mathcal{A}}).
\]
It may be more intuitive to reason equivalently as:
\[
\sum_{o\in\mathrm{OPT}}\operatorname{gain}_{\tau(o)}(o)
+f(\mathcal{S}_{\mathcal{A}})\geq f(\mathrm{OPT}).
\]
Some may see immediately why this inequality holds, if a village covered by
an optimal pair $o$ is not counted in any of the $\operatorname{gain}_{\tau(o)}(o)$ because
greedy has already removed it from the uncovered villages, then it is already
counted in $f(\mathcal{S}_{\mathcal{A}})$.

To put it formally, let $N_{\mathcal{A}}$ be the villages covered by greedy and
$N_{\tau(o)}$ the uncovered villages immediately before the selection that
removes $o$ (or at termination if $o$ remains):
\[
N_{\mathcal{A}}:=\bigcup_{t\in\mathcal{S}_{\mathcal{A}}}V[t],\qquad
N_{\mathrm{OPT}}:=\bigcup_{o\in\mathrm{OPT}}V[o],\qquad
N_R:=\bigcup_{o\in\mathrm{OPT}}\bigl(V[o]\cap N_{\tau(o)}\bigr).
\]
$N_{\mathrm{OPT}}$ is the set of villages covered by the optimal solution.
$N_R$ is the set of villages covered by an optimal pair that are still uncovered
by greedy when that pair is removed, namely the villages counted in the gain terms.
\[
N_{\mathrm{OPT}}\subseteq N_{\mathcal{A}}\cup N_R.
\]
\[
\begin{aligned}
\sum_{o\in\mathrm{OPT}}\operatorname{gain}_{\tau(o)}(o)
+f(\mathcal{S}_{\mathcal{A}})
&\geq\sum_{v\in N_R}P[v]+\sum_{v\in N_{\mathcal{A}}}P[v]\\
&\geq\sum_{v\in N_{\mathcal{A}}\cup N_R}P[v]\\
&\geq\sum_{v\in N_{\mathrm{OPT}}}P[v]
=f(\mathrm{OPT}).
\end{aligned}
\]
Thus, we have proved:
\[
2f(\mathcal{S}_{\mathcal{A}})
\geq\sum_{o\in\mathrm{OPT}}\operatorname{gain}_{\tau(o)}(o)
\geq f(\mathrm{OPT})-f(\mathcal{S}_{\mathcal{A}}).
\]
Rearranging gives:
\[
3f(\mathcal{S}_{\mathcal{A}})\geq f(\mathrm{OPT})
\quad\Longrightarrow\quad
f(\mathcal{S}_{\mathcal{A}})\geq\frac{f(\mathrm{OPT})}{3}.
\]
So our algorithm $\mathcal{A}$ is a $3$-approximation.\hfill$\square$
\end{enumerate}

\end{document}
